% Source: https://artofproblemsolving.com/community/c3h2524678_prime_numbers_pattern
% Stable identity: aops-topic-2524678-q01
\textbf{(a)} Since $10a\equiv-1\pmod n$,
\[10(Q-aR)=10Q-10aR\equiv10Q+R=N\pmod n.\]
So $n\mid N$ if and only if $n\mid10(Q-aR)$, and because $\gcd(n,10)=1$ this happens if and only if $n\mid Q-aR$. (This is the familiar kind of divisibility test: for $7$, $a=2$ since $10\cdot2=20\equiv-1\pmod7$, so $7\mid10Q+R$ exactly when $7\mid Q-2R$.)

\textbf{(b)} Let $n=10q+r$ be a two-digit number with $\gcd(n,10)=1$. Since $10$ is invertible modulo $n$, there is an $a$ with $10a\equiv-1\pmod n$, and we may take $1\le a\le n-1$ (as $a\equiv0$ would give $0\equiv-1$). Put $b=n-a\ge1$. Applying (a) to $N=n$ itself (so $Q=q$, $R=r$) gives $n\mid q-ar$, and
\[q+br=q+nr-ar\equiv q-ar\equiv0\pmod n.\]
So $n$ is nice. For example $13$: $10\cdot9=90\equiv-1\pmod{13}$, giving the pair $9,4$; for $17$: $10\cdot5=50\equiv-1\pmod{17}$, giving $5,12$. These are exactly the pairs found by the student.

The property depends only on $\gcd(n,10)=1$, not on primality. For the composite number $21$: $10\cdot2=20\equiv-1\pmod{21}$, so take $a=2$, $b=19$: $2-2\cdot1=0$ and $2+19\cdot1=21$ are both divisible by $21$. So
\[\boxed{21=3\cdot7\text{ is nice (with }a=2,\ b=19\text{), and the conjecture is false.}}\]
Every other composite two-digit number coprime to $10$ (such as $27$, $33$, $39$, $49$, $\dots$, $99$) is a counterexample too.
